% GENERATED FILE. Edit canonical lessons, then run npm run generate:books.
% canonical-source-hash: fnv1a64:be7c931c5d1b6e86
% canonical-lessons: GU-C140-kapvu, GU-C140-sivvu, GU-C140-bharvu, GU-C140-khecvu, GU-C140-phekvu

\chapter{To cut, To sew, To fill, To pull, To throw}
\label{ch:gu-to-cut-to-sew-to-fill-to-pull-to-throw}
\hlchaptermodality{140}

\begin{chapteropening}
\textbf{By the end of this chapter:} \emph{I can say to cut, to sew, to fill, to pull and to throw.}

Complete the last lesson of chapter 140: I can say to cut, to sew, to fill, to pull and to throw.
\end{chapteropening}

\section[\texorpdfstring{\gu{કાપવું} (kāpvũ)}{કાપવું (kāpvũ)}]{\gu{કાપવું} (kāpvũ) — to cut}

\label{lesson:GU-C140-kapvu}



\begin{quote}
Before the new one: say the Gujarati for to use, then the Gujarati for to make.
\end{quote}

\subsection*{You'll want to know: \gu{કાપવું}}
\textbf{\gu{કાપવું}} (\emph{kāpvũ}) — \textquotedblleft{}to cut\textquotedblright{}.

\begin{cousinweb}
\textbf{\gu{કાપવું}} (\emph{kāpvũ}) is a tadbhava, a word inherited from Sanskrit and reshaped by sound change. It comes from \emph{kalpayati}, \textquotedblleft{}sets in order; trims, cuts\textquotedblright{} (CDIAL 2944).
\end{cousinweb}

\begin{grammarlens}[title={What you've built}]
One more word for this chapter.
\end{grammarlens}

\subsection*{Your turn}
\begin{itemize}
\raggedright
  \item \emph{Say it:} \emph{kāpvũ}
  \item \emph{Say it:} \emph{kāpvũ}, once more
  \item \emph{Say it:} say \emph{kāpvũ}
  \item \emph{Recall:} say the Gujarati for to use, then the Gujarati for to make, then say \emph{kāpvũ} again
\end{itemize}

\begin{tcolorbox}[breakable,colback=teal!4,colframe=teal!35!black,title={Before you move on}]
What does \gu{કાપવું} mean? (\textquotedblleft{}To cut\textquotedblright{}.) Say it once more.
\end{tcolorbox}
\section[\texorpdfstring{\gu{સીવવું} (sīvvũ)}{સીવવું (sīvvũ)}]{\gu{સીવવું} (sīvvũ) — to sew}

\label{lesson:GU-C140-sivvu}



\begin{quote}
Before the new one: say the Gujarati for to make, then the Gujarati for to cut.
\end{quote}

\subsection*{You'll want to know: \gu{સીવવું}}
\textbf{\gu{સીવવું}} (\emph{sīvvũ}) — \textquotedblleft{}to sew\textquotedblright{}.

\begin{cousinweb}
\textbf{\gu{સીવવું}} (\emph{sīvvũ}) is a tadbhava. CDIAL lists it under Sanskrit \emph{sīvyati}, \textquotedblleft{}sews\textquotedblright{}, noting that it may come instead from the form \emph{sīvayati} (CDIAL 13444).
\end{cousinweb}

\begin{grammarlens}[title={What you've built}]
One more word for this chapter.
\end{grammarlens}

\subsection*{Your turn}
\begin{itemize}
\raggedright
  \item \emph{Say it:} \emph{sīvvũ}
  \item \emph{Say it:} \emph{sīvvũ}, once more
  \item \emph{Say it:} say \emph{sīvvũ}
  \item \emph{Recall:} say the Gujarati for to make, then the Gujarati for to cut, then say \emph{sīvvũ} again
\end{itemize}

\begin{tcolorbox}[breakable,colback=teal!4,colframe=teal!35!black,title={Before you move on}]
What does \gu{સીવવું} mean? (\textquotedblleft{}To sew\textquotedblright{}.) Say it once more.
\end{tcolorbox}
\section[\texorpdfstring{\gu{ભરવું} (bharvũ)}{ભરવું (bharvũ)}]{\gu{ભરવું} (bharvũ) — to fill}

\label{lesson:GU-C140-bharvu}



\begin{quote}
Before the new one: say the Gujarati for to cut, then the Gujarati for to sew.
\end{quote}

\subsection*{You'll want to know: \gu{ભરવું}}
\textbf{\gu{ભરવું}} (\emph{bharvũ}) — \textquotedblleft{}to fill\textquotedblright{}.

\begin{cousinweb}
\textbf{\gu{ભરવું}} (\emph{bharvũ}) is a tadbhava, a word inherited from Sanskrit and reshaped by sound change. It comes from \emph{bharati}, \textquotedblleft{}bears, brings; fills\textquotedblright{} (CDIAL 9397).
\end{cousinweb}

\begin{grammarlens}[title={What you've built}]
One more word for this chapter.
\end{grammarlens}

\subsection*{Your turn}
\begin{itemize}
\raggedright
  \item \emph{Say it:} \emph{bharvũ}
  \item \emph{Say it:} \emph{bharvũ}, once more
  \item \emph{Say it:} say \emph{bharvũ}
  \item \emph{Recall:} say the Gujarati for to cut, then the Gujarati for to sew, then say \emph{bharvũ} again
\end{itemize}

\begin{tcolorbox}[breakable,colback=teal!4,colframe=teal!35!black,title={Before you move on}]
What does \gu{ભરવું} mean? (\textquotedblleft{}To fill\textquotedblright{}.) Say it once more.
\end{tcolorbox}
\section[\texorpdfstring{\gu{ખેંચવું} (khẽcvũ)}{ખેંચવું (khẽcvũ)}]{\gu{ખેંચવું} (khẽcvũ) — to pull}

\label{lesson:GU-C140-khecvu}



\begin{quote}
Before the new one: say the Gujarati for to sew, then the Gujarati for to fill.
\end{quote}

\subsection*{You'll want to know: \gu{ખેંચવું}}
\textbf{\gu{ખેંચવું}} (\emph{khẽcvũ}) — \textquotedblleft{}to pull\textquotedblright{}.

\begin{cousinweb}
CDIAL lists Gujarati \emph{khīṁcvũ}, \textquotedblleft{}to pull\textquotedblright{}, under a reconstructed \emph{*khiñc}, \textquotedblleft{}drag, pull\textquotedblright{}, and the noun \emph{khẽc}, \textquotedblleft{}pulling\textquotedblright{}, under its variant \emph{*khaiñc} (CDIAL 3881). \textbf{\gu{ખેંચવું}} (\emph{khẽcvũ}) belongs with them.
\end{cousinweb}

\begin{grammarlens}[title={What you've built}]
One more word for this chapter.
\end{grammarlens}

\subsection*{Your turn}
\begin{itemize}
\raggedright
  \item \emph{Say it:} \emph{khẽcvũ}
  \item \emph{Say it:} \emph{khẽcvũ}, once more
  \item \emph{Say it:} say \emph{khẽcvũ}
  \item \emph{Recall:} say the Gujarati for to sew, then the Gujarati for to fill, then say \emph{khẽcvũ} again
\end{itemize}

\begin{tcolorbox}[breakable,colback=teal!4,colframe=teal!35!black,title={Before you move on}]
What does \gu{ખેંચવું} mean? (\textquotedblleft{}To pull\textquotedblright{}.) Say it once more.
\end{tcolorbox}
\section[\texorpdfstring{\gu{ફેંકવું} (phẽkvũ)}{ફેંકવું (phẽkvũ)}]{\gu{ફેંકવું} (phẽkvũ) — to throw}

\label{lesson:GU-C140-phekvu}



\begin{quote}
Before the new one: say the Gujarati for to fill, then the Gujarati for to pull.
\end{quote}

\subsection*{You'll want to know: \gu{ફેંકવું}}
\textbf{\gu{ફેંકવું}} (\emph{phẽkvũ}) — \textquotedblleft{}to throw\textquotedblright{}.

\begin{cousinweb}
\textbf{\gu{ફેંકવું}} (\emph{phẽkvũ}) is a tadbhava. CDIAL lists it under \emph{*pheṅk-}, one of the forms it reconstructs under \emph{*phekk}, \textquotedblleft{}throw\textquotedblright{} (CDIAL 9106).
\end{cousinweb}

\begin{grammarlens}[title={What you've built}]
One more word for this chapter.
\end{grammarlens}

\subsection*{Your turn}
\begin{itemize}
\raggedright
  \item \emph{Say it:} \emph{phẽkvũ}
  \item \emph{Say it:} \emph{phẽkvũ}, once more
  \item \emph{Say it:} say \emph{phẽkvũ}
  \item \emph{Recall:} say the Gujarati for to fill, then the Gujarati for to pull, then say \emph{phẽkvũ} again
\end{itemize}

\begin{tcolorbox}[breakable,colback=teal!4,colframe=teal!35!black,title={Before you move on}]
What does \gu{ફેંકવું} mean? (\textquotedblleft{}To throw\textquotedblright{}.) Say it once more.
\end{tcolorbox}
